\documentclass[12pt, a4paper]{article}
\usepackage[margin=2.5cm]{geometry}
\usepackage{booktabs}
\usepackage{enumitem}

\title{\textbf{R25 Sprint Burndown --- Logic}}
\author{DXC Service Desk Analytics}
\date{Sprint R25 \quad|\quad 08 March -- 04 April 2026}

\begin{document}
\maketitle

% ─────────────────────────────────────────────────────────────────────────────
\section{Scope}

The burndown only looks at the \textbf{42 tickets} committed in the
Release Lifecycle file (sheet \emph{Scope UPDATE}). No other tickets are included.

% ─────────────────────────────────────────────────────────────────────────────
\section{Story Points per Ticket}

Story points are read from the \textbf{Jira extract} first.
If a ticket has no value there, the \textbf{planning file} value is used instead.

% ─────────────────────────────────────────────────────────────────────────────
\section{When is a Ticket Considered Done?}

Three date columns are checked in order, taking the first one that is filled:

\begin{enumerate}
  \item \texttt{qualification\_date} --- ticket passed the QA gate
  \item \texttt{resolved} --- Jira resolution date
  \item \texttt{closure\_date} --- formally closed in the workflow
\end{enumerate}

That date is then accepted or rejected based on this rule:

\begin{itemize}
  \item Date is \textbf{within the sprint} (08 Mar -- 04 Apr) $\rightarrow$ accepted.
  \item Date is \textbf{outside the sprint} $\rightarrow$ only accepted if the
        ticket status is \emph{Closed}, \emph{Cancelled}, \emph{Done}, or \emph{Resolved}.
  \item Anything else $\rightarrow$ ignored.
\end{itemize}

\textbf{Why?} Some tickets have a \texttt{qualification\_date} left over
from a previous sprint, even though they were re-opened and are still in
progress. Without this check, those tickets would wrongly count as already
delivered from day one.

% ─────────────────────────────────────────────────────────────────────────────
\section{What is Computed Each Day}

For every calendar day in the sprint, the system counts all tickets whose
validated completion date falls on or before that day (by end of day).
From that count the following columns are derived:

\medskip
\begin{center}
\begin{tabular}{ll}
\toprule
\textbf{Column} & \textbf{How it is calculated} \\
\midrule
\texttt{resolved\_pts}           & Sum of story points of completed tickets \\
\texttt{remaining\_pts}          & Total SP minus resolved SP \\
\texttt{ideal\_pts}              & Total SP spread evenly across sprint days (linear decrease to 0) \\
\texttt{resolved\_tickets}       & Number of completed tickets \\
\texttt{remaining\_tickets}      & 42 minus resolved tickets \\
\texttt{ideal\_tickets}          & 42 spread evenly across sprint days (linear decrease to 0) \\
\bottomrule
\end{tabular}
\end{center}

\medskip
The \textbf{ideal} columns represent what perfect, steady progress would look
like --- starting at the full scope on day one and reaching zero on the last day.

\end{document}
